\ifdefined\gkver\else\def\gkver{student}\fi
\documentclass[\gkver]{gaokaozhenti}
\newcommand{\ccnum}[1]{{\fontspec{Noto Sans CJK JP}#1}}
\begin{document}

\chapter{2018年全国理综Ⅰ}
%% paperType: 理综物理部分

\begin{enumerate}
\item
%% number: 14
%% paperName: 2018年全国理综Ⅰ第 14 题 6 分
%% typeId: 1
%% type: 单选题
%% chapter: 功和能
%% point: 动能定理
%% method: 无
%% score: 6
%% degree: 750
%% duplicateId: 0
%% body:
高铁列车在启动阶段的运动可看作初速度为零的匀加速直线运动，在启动阶段列车的动能\xzanswer{B}

\fourchoices[answer=B]
{与它所经历的时间成正比}
{与它的位移成正比}
{与它的速度成正比}
{与它的动量成正比}

%% answer: B
%% memo:
\memoanswer{
列车做匀加速直线运动，由 $v = at$ 可知列车的动能 $E_k = \frac{1}{2}mv^2 = \frac{1}{2}ma^2t^2 \propto t^2$，故 A 错误；由 $v^2 = 2ax$ 可知列车的动能 $E_k = \frac{1}{2}mv^2 = max \propto x$，故 B 正确；列车的动能 $E_k = \frac{1}{2}mv^2 \propto v^2$，故 C 错误；列车的动能 $E_k = \frac{1}{2}mv^2 = \frac{p^2}{2m} \propto p^2$，故 D 错误。
}

\item
%% number: 15
%% paperName: 2018年全国理综Ⅰ第 15 题 6 分
%% typeId: 1
%% type: 单选题
%% chapter: 运动和力
%% point: 牛顿定律的应用
%% method: 无
%% score: 6
%% degree: 550
%% duplicateId: 0
%% body:
如图，轻弹簧的下端固定在水平桌面上，上端放有物块 P，系统处于静止状态，现用一竖直向上的力 $F$ 作用在 P 上，使其向上做匀加速直线运动，以 $x$ 表示 P 离开静止位置的位移，在弹簧恢复原长前，下列表示 $F$ 和 $x$ 之间关系的图像可能正确的是\xzanswer{A}

\onepicture[w=2.4cm, align=right]{figs/15.png}

\fourchoices[answer=A, ispicture=true, h=1.7cm]
{figs/15a.png}
{figs/15b.png}
{figs/15c.png}
{figs/15d.png}

%% answer: A
%% memo:
\memoanswer{
设物块静止时弹簧的形变量为 $x_0$，则有 $mg = kx_0$。物块做匀加速直线运动，由牛顿第二定律得 $F - mg + k(x_0 - x) = ma$，解得 $F = ma + kx$，所以 $F-x$ 图线是不过原点的倾斜直线，故 A 正确，BCD 错误。
}

\item
%% number: 16
%% paperName: 2018年全国理综Ⅰ第 16 题 6 分
%% typeId: 1
%% type: 单选题
%% chapter: 电场
%% point: 库仑定律
%% method: 无
%% score: 6
%% degree: 550
%% duplicateId: 0
%% body:
如图，三个固定的带电小球 a、b 和 c，相互间的距离分别为 $ab = 5\Ucm$，$bc = 3\Ucm$，$ca = 4\Ucm$。小球 c 所受库仑力的合力的方向平行于 a、b 的连线。设小球 a、b 所带电荷量的比值的绝对值为 $k$，则\xzanswer{D}

\onepicture[w=6.5cm]{\includesvg{figs/16}}

\fourchoices[answer=D]
{a、b 的电荷同号，$k = \frac{16}{9}$}
{a、b 的电荷异号，$k = \frac{16}{9}$}
{a、b 的电荷同号，$k = \frac{64}{27}$}
{a、b 的电荷异号，$k = \frac{64}{27}$}

%% answer: D
%% memo:
\memoanswer{
由题意可知小球 c 受到的库仑力的合力可能水平向左，也可能水平向右，受力如图所示。

\onepicture[w=4.2cm]{figs/16_an.jpg}

若合力水平向左，可知小球 a 对小球 c 吸引，小球 b 对小球 c 排斥，所以 a、b 带异种电荷；由平衡条件可知 $\tan 37^\circ = \frac{F_1}{F_2}$，由库仑力公式有 $F_1 = \frac{k' q_b q_c}{(bc)^2}$，$F_2 = \frac{k' q_a q_c}{(ac)^2}$，联立解得 $\frac{q_a}{q_b} = \frac{64}{27}$；同理，若库仑力的合力水平向右，可知小球 a 对小球 c 排斥，小球 b 对小球 c 吸引，所以 a、b 带异种电荷，同样有 $\tan 37^\circ = \frac{F_1'}{F_2'}$，仍可推得 $\frac{q_a}{q_b} = \frac{64}{27}$，故 D 正确，ABC 错误。
}

\item
%% number: 17
%% paperName: 2018年全国理综Ⅰ第 17 题 6 分
%% typeId: 1
%% type: 单选题
%% chapter: 电磁感应
%% point: 电磁感应的电量
%% method: 无
%% score: 6
%% degree: 950
%% duplicateId: 0
%% body:
如图，导体轨道 OPQS 固定，其中 PQS 是半圆弧，Q 为半圆弧的中心，O 为圆心。轨道的电阻忽略不计。OM 是有一定电阻、可绕 O 转动的金属杆。M 端位于 PQS 上，OM 与轨道接触良好。空间存在与半圆所在平面垂直的匀强磁场，磁感应强度的大小为 $B$，现使 OQ 位置以恒定的角速度逆时针转到 OS 位置并固定（过程Ⅰ）；再使磁感应强度的大小以一定的变化率从 $B$ 增加到 $B'$（过程Ⅱ）。在过程Ⅰ、Ⅱ中，流过 OM 的电荷量相等，则 $\frac{B'}{B}$ 等于\xzanswer{B}

\onepicture{figs/17.png}

\fourchoices[answer=B]
{$\frac{5}{4}$}
{$\frac{3}{2}$}
{$\frac{7}{4}$}
{2}

%% answer: B
%% memo:
\memoanswer{
OM 从 OQ 位置以恒定的角速度逆时针转动到 OS 位置，通过 OM 的电荷量 $q_1 = \frac{\overline{E}}{R}\Delta t = \frac{\Delta\phi}{\Delta t}\Delta t = \frac{\Delta\phi}{R} = \frac{B \cdot \frac{\pi r^2}{4}}{R}$；磁感应强度的大小以一定的变化率从 $B$ 增加到 $B'$ 的过程中，通过 OM 的电荷量 $q_2 = \frac{\Delta\phi}{R} = \frac{S\Delta B}{R} = \frac{\frac{\pi}{2}r^2(B'-B)}{R}$，并且 $q_1 = q_2$，联立解得 $\frac{B'}{B} = \frac{3}{2}$，故 B 正确。
}

\item
%% number: 18
%% paperName: 2018年全国理综Ⅰ第 18 题 6 分
%% typeId: 1
%% type: 单选题
%% chapter: 功和能
%% point: 动能定理
%% method: 无
%% score: 6
%% degree: 450
%% duplicateId: 0
%% body:
如图，abc 是竖直面内的光滑固定轨道，ab 水平，长度为 $2R$，bc 是半径为 $R$ 的四分之一的圆弧，与 ab 相切于 b 点。一质量为 $m$ 的小球始终受到与重力大小相等的水平外力的作用，自 a 点处从静止开始向右运动，重力加速度大小为 $g$。小球从 a 点开始运动到其轨迹最高点，机械能的增量为\xzanswer{C}

\onepicture{figs/18.png}

\fourchoices[answer=C]
{$2mgR$}
{$4mgR$}
{$5mgR$}
{$6mgR$}

%% answer: C
%% memo:
\memoanswer{
小球从 a 到 c 的过程，由动能定理得 $F\cdot3R - mgR = \frac{1}{2}mv_c^2 - 0$，由题意知 $F$ 的大小为 $mg$，解得 $v_c = 2\sqrt{gR}$。小球离开 c 点后在竖直方向做竖直上抛运动，水平方向做匀加速直线运动，在竖直方向有 $v_c = gt$，在水平方向上有 $x = \frac{1}{2}at^2 = \frac{Ft^2}{2m}$，联立解得 $x = 2R$，可知小球机械能的增量为 $\Delta E = F(3R + x) = 5mgR$，故 C 正确。
}

\item
%% number: 19
%% paperName: 2018年全国理综Ⅰ第 19 题 6 分
%% typeId: 2
%% type: 多选题
%% chapter: 电磁感应
%% point: 楞次定律
%% method: 无
%% score: 6
%% degree: 550
%% duplicateId: 0
%% body:
如图，两个线圈绕在同一根铁芯上，其中一线圈通过开关与电源连接，另一线圈与远处沿南北方向水平放置在纸面内的直导线连接成回路。将一小磁针悬挂在直导线正上方，开关未闭合时小磁针处于静止状态。下列说法正确的是\xzanswer{AD}

\onepicture{figs/19.png}

\fourchoices[answer=AD]
{开关闭合后的瞬间，小磁针的 N 极朝垂直纸面向里的方向转动}
{开关闭合并保持一段时间后，小磁针的 N 极指向垂直纸面向里的方向}
{开关闭合并保持一段时间后，小磁针的 N 极指向垂直纸面向外的方向}
{开关闭合并保持一段时间再断开后的瞬间，小磁针的 N 极朝垂直纸面向外的方向转动}

%% answer: AD
%% memo:
\memoanswer{
开关闭合的瞬间，左侧的线圈中磁通量变化，产生感应电动势和感应电流，由楞次定律可判断出直导线中电流方向为由南向北，由安培定则可判断出小磁针处的磁场方向垂直纸面向里，小磁针的 N 极朝垂直纸面向里的方向转动，选项 A 正确；

开关闭合并保持一段时间后，左侧线圈中磁通量不变，线圈中感应电动势和感应电流为零，直导线中电流为零，小磁针恢复到原来状态，选项 BC 错误；

开关闭合并保持一段时间再断开后的瞬间，左侧的线圈中磁通量变化，产生感应电动势和感应电流，由楞次定律可判断出直导线中电流方向为由北向南，由安培定则可判断出小磁针处的磁场方向垂直纸面向外，小磁针的 N 极朝垂直纸面向外的方向转动，选项 D 正确。
}

\item
%% number: 20
%% paperName: 2018年全国理综Ⅰ第 20 题 6 分
%% typeId: 2
%% type: 多选题
%% chapter: 曲线运动
%% point: 双星问题
%% method: 无
%% score: 6
%% degree: 450
%% duplicateId: 0
%% body:
2017 年，人类第一次直接探测到来自双中子星合并的引力波。根据科学家们复原的过程，在两颗中子星合并前约 $100\Us$ 时，它们相距约 $400\Ukm$，绕二者连线上的某点每秒转动 12 圈，将两颗中子星都看作是质量均匀分布的球体，由这些数据、万有引力常量并利用牛顿力学知识，可以估算出这一时刻两颗中子星\xzanswer{BC}

\fourchoices[answer=BC]
{质量之积}
{质量之和}
{速率之和}
{各自的自转角速度}

%% answer: BC
%% memo:
\memoanswer{
由题意知两中子星的距离 $L$ 和周期 $T$，由万有引力提供向心力得 $\frac{Gm_1m_2}{L^2} = m_1\frac{4\pi^2}{T^2}r_1$，$\frac{Gm_1m_2}{L^2} = m_2\frac{4\pi^2}{T^2}r_2$，其中 $r_1 + r_2 = L$，联立解得 $\frac{G(m_1+m_2)}{L^2} = \frac{4\pi^2 L}{T^2}$，可以求出它们的质量之和 $m_1 + m_2 = \frac{4\pi^2 L^3}{GT^2}$，故 B 正确；由 $v_1 = \frac{2\pi r_1}{T}$，$v_2 = \frac{2\pi r_2}{T}$，可得它们的速率之和 $v_1 + v_2 = \frac{2\pi L}{T}$，故 C 正确；由于 $r_1$、$r_2$ 的具体大小未知，无法求出它们的质量之积 $m_1m_2 = \frac{16\pi^4 L^4 r_1 r_2}{G^2T^4}$，故 A 错误；题中条件与中子星自转无关，无法求出自转角速度，故 D 错误。
}

\item
%% number: 21
%% paperName: 2018年全国理综Ⅰ第 21 题 6 分
%% typeId: 2
%% type: 多选题
%% chapter: 电场
%% point: 电势能电势差电场力做功
%% method: 无
%% score: 6
%% degree: 350
%% duplicateId: 0
%% body:
图中虚线 a、b、c、d、f 代表匀强电场内间距相等的一组等势面，已知平面 b 上的电势为 $2\UV$。一电子经过 a 时的动能为 $10\UeV$，从 a 到 d 的过程中克服电场力所做的功为 $6\UeV$。下列说法正确的是\xzanswer{AB}

\onepicture[w=6.0cm]{\includesvg{figs/21}}

\fourchoices[answer=AB]
{平面 c 上的电势为零}
{该电子可能到达不了平面 f}
{该电子经过平面 d 时，其电势能为 $4\UeV$}
{该电子经过平面 b 时的速率是经过 d 时的 2 倍}

%% answer: AB
%% memo:
\memoanswer{
由题可知 $U_{ad} = \frac{W_{ad}}{q} = \frac{-6\UeV}{-e} = 6\UV$，由图可知 $U_{ad} = 3U_{ab} = 3U_{bc}$，所以 $\varphi_a - \varphi_b = \varphi_b - \varphi_c = 2\UV$，故平面 c 上的电势 $\varphi_c = 0$，故 A 正确；

电子进入电场的方向不确定，可能在到达平面 f 之前沿电场方向的速度减为零，垂直于电场方向的速度不为零，故 B 正确；

由以上分析可得平面 d 的电势 $\varphi_d = -2\UV$，所以电子在平面 d 的电势能 $E_p = -e\varphi_d = 2\UeV$，故 C 错误；电子从 a 到 b 由动能定理可得 $W_{ab} = E_{kb} - E_{ka}$，解得 $E_{kb} = 8\UeV$；电子从 a 到 d 由动能定理可得 $W_{ad} = -U_{ad}e = E_{kd} - E_{ka}$，解得 $E_{kd} = 4\UeV$，则 $v_b = \sqrt{2}v_d$，故 D 错误。
}

\item
%% number: 22
%% paperName: 2018年全国理综Ⅰ第 22 题 5 分
%% typeId: 4
%% type: 实验题
%% chapter: 力物体的平衡
%% point: 胡克定律
%% method: 无
%% score: 5
%% degree: 850
%% duplicateId: 0
%% body:
如图（a），一弹簧上端固定在支架顶端，下端悬挂一托盘；一标尺由游标和主尺构成，主尺竖直固定在弹簧左边；托盘上方固定有一能与游标刻度线准确对齐的装置，简化为图中的指针。现要测量图（a）中弹簧的劲度系数，当托盘内没有砝码时，移动游标，使其零刻度线对准指针，此时标尺读数为 $1.950\Ucm$；当托盘内放有质量为 $0.100\Ukg$ 的砝码时，移动游标，再次使其零刻度线对准指针，标尺示数如图（b）所示，其读数为\tkanswer[2.0cm]{$3.775$} $\Ucm$。当地的重力加速度大小为 $9.80\Umsq$，此弹簧的劲度系数为\tkanswer[2.0cm]{$53.7$} $\UNm$（保留 3 位有效数字）。

\twopicture[labelA=2018全国理综Ⅰ22a, labelB=2018全国理综Ⅰ22b, numA={（a）}, numB={（b）}, hA=5.5cm, hB=5.5cm, align=right]
{figs/22a.png}
{figs/22b.png}

%% answer:
%% memo:
\memoanswer{
主尺读数为 $3.7\Ucm$，游标上第 15 个刻度和主尺上某一刻度对齐，所以游标读数为 $15 \times 0.05\Umm = 0.75\Umm = 0.075\Ucm$，所以最终读数为 $3.7\Ucm + 0.075\Ucm = 3.775\Ucm$；由胡克定律得 $k = \frac{F}{\Delta x} = \frac{mg}{L_1 - L_0} = \frac{0.100 \times 9.8}{(3.775 - 1.950) \times 10^{-2}}\UNm \approx 53.7\UNm$。
}

\item
%% number: 23
%% paperName: 2018年全国理综Ⅰ第 23 题 10 分
%% typeId: 4
%% type: 实验题
%% chapter: 电路
%% point: 伏安法测电阻
%% method: 无
%% score: 10
%% degree: 550
%% duplicateId: 0
%% body:
某实验小组利用如图（a）所示的电路探究在 $25\UdC\sim80\UdC$ 范围内某热敏电阻的温度特性，所用器材有：置于温控室（图中虚线区域）中的热敏电阻 $R_T$，其标称值（$25\UdC$ 时的阻值）为 $900.0\UO$；电源 $E$（$6\UV$，内阻可忽略）；电压表 V（量程 $150\UmV$）；定值电阻 $R_0$（阻值 $20.0\UO$），滑动变阻器 $R_1$（最大阻值为 $1000\UO$）；电阻箱 $R_2$（阻值范围 $0\sim999.9\UO$）；单刀开关 $S_1$，单刀双掷开关 $S_2$。

实验时，先按图（a）连接好电路，再将温控室的温度 $t$ 升至 $80.0\UdC$，将 $S_2$ 与 1 端接通，闭合 $S_1$，调节 $R_1$ 的滑片位置，使电压表读数为某一值 $U_0$；保持 $R_1$ 的滑片位置不变，将 $R_2$ 置于最大值，将 $S_2$ 与 2 端接通，调节 $R_2$，使电压表读数仍为 $U_0$；断开 $S_1$，记下此时 $R_2$ 的读数，逐步降低温控室的温度 $t$，得到相应温度下 $R_2$ 的阻值，直至温度降到 $25.0\UdC$，实验得到的 $R_2-t$ 数据见下表。

\begin{table}[htbp]
\centering
\begin{tblr}{|c|c|c|c|c|c|c|c|}
\hline
$t/\UdC[a]$ & 25.0 & 30.0 & 40.0 & 50.0 & 60.0 & 70.0 & 80.0 \\
\hline
$R_2/\UO[a]$ & 900.0 & 680.0 & 500.0 & 390.0 & 320.0 & 270.0 & 240.0 \\
\hline
\end{tblr}
\end{table}

回答下列问题：
\begin{enumerate}
	\item
	在闭合 $S_1$ 前，图（a）中 $R_1$ 的滑片应移动到\tkanswer[1.4cm]{b}（填“a”或“b”）端；
	\item
	在图（b）的坐标纸上补齐数据表中所给数据点，并做出 $R_2-t$ 曲线；
	\item
	由图（b）可得到 $R_T$ 在 $25\UdC\sim80\UdC$ 范围内的温度特性，当 $t = 44.0\UdC$ 时，可得 $R_T = $\tkanswer[1.6cm]{$450$} $\UO$；
	\item
	将 $R_T$ 握于手心，手心温度下 $R_2$ 的相应读数如图（c）所示，该读数为\tkanswer[1.8cm]{$620.0$} $\UO$，则手心温度为\tkanswer[1.6cm]{$33.0$} $\UdC$。
\end{enumerate}

\threepicture[labelA=2018全国理综Ⅰ23a, labelB=2018全国理综Ⅰ23b, labelC=2018全国理综Ⅰ23c, numA={（a）}, numB={（b）}, numC={（c）}, hA=3.2cm, hB=3.2cm, hC=3.2cm, align=right]
{figs/23a.png}
{\includesvg{figs/23b}}
{figs/23c.png}

%% answer:
\jdanswer{
\begin{enumerate}
	\item
	b

	\item
	如图所示

	\item
	$450$

	\item
	$620.0$，$33.0$

\end{enumerate}
}
%% memo:
\memoanswer{
（1）滑动变阻器采用限流式接法，在闭合开关前，从保护电路的角度考虑，滑动变阻器接入电路的阻值应调到最大值，故 $R_1$ 的滑片应移动到 b 端。

（2）先描点，之后在坐标纸上用平滑的曲线连接各点，如图所示。

\onepicture[w=6.5cm]{figs/23_an.png}

（3）由图（b）可知，当 $t = 44.0\UdC$ 时，$R_T = R_2 = 450\UO$。

（4）由图（c）可知，电阻箱示数为 $R_T = 6\times100\UO + 2\times10\UO + 0\times1\UO + 0\times0.1\UO = 620.0\UO$，由图（b）可得当 $R_T = R_2 = 620.0\UO$ 时，手心的温度 $t = 33.0\UdC$。
}

\item
%% number: 24
%% paperName: 2018年全国理综Ⅰ第 24 题 12 分
%% typeId: 6
%% type: 计算题
%% chapter: 运动和力
%% point: 动量和能量
%% method: 无
%% score: 12
%% degree: 550
%% duplicateId: 0
%% body:
一质量为 $m$ 的烟花弹获得动能 $E$ 后，从地面竖直升空，当烟花弹上升的速度为零时，弹中火药爆炸将烟花弹炸为质量相等的两部分，两部分获得的动能之和也为 $E$，且均沿竖直方向运动。爆炸时间极短，重力加速度大小为 $g$，不计空气阻力和火药的质量，求：
\begin{enumerate}
	\item
	烟花弹从地面开始上升到弹中火药爆炸所经过的时间；
	\item
	爆炸后烟花弹向上运动的部分距地面的最大高度。
\end{enumerate}

%% answer:
\jdanswer{
\begin{enumerate}
	\item
	$t = \frac{1}{g}\sqrt{\frac{2E}{m}}$

	\item
	$h = \frac{2E}{mg}$

\end{enumerate}
}
%% memo:
\memoanswer{
（1）设烟花弹上升的初速度为 $v_0$，由题给条件得

$E=\frac{1}{2}mv_0^2$\ccnum{①}，

设烟花弹从地面开始上升到火药爆炸所用的时间为 $t$，由运动学公式得

$0-v_0=-gt$\ccnum{②}，

联立\ccnum{①}\ccnum{②}式得 $t=\frac{1}{g}\sqrt{\frac{2E}{m}}$\ccnum{③}。

（2）设爆炸时烟花弹距地面的高度为 $h_1$，由机械能守恒定律得

$E=mgh_1$\ccnum{④}，

火药爆炸后，烟花弹上、下两部分均沿竖直方向运动，设爆炸后瞬间其速度分别为 $v_1$ 和 $v_2$，由题给条件和动量守恒得

$\frac{1}{4}mv_1^2+\frac{1}{4}mv_2^2=E$\ccnum{⑤}，

$\frac{1}{2}mv_1+\frac{1}{2}mv_2=0$\ccnum{⑥}，

由\ccnum{⑥}式知，烟花弹两部分的速度方向相反，向上运动部分做竖直上抛运动。设爆炸后烟花弹上部分继续上升的高度为 $h_2$，由机械能守恒定律得

$\frac{1}{4}mv_1^2=\frac{1}{2}mgh_2$\ccnum{⑦}，

联立\ccnum{④}\ccnum{⑤}\ccnum{⑥}\ccnum{⑦}式得，烟花弹上部分距地面的最大高度为

$h=h_1+h_2=\frac{2E}{mg}$\ccnum{⑧}。
}

\item
%% number: 25
%% paperName: 2018年全国理综Ⅰ第 25 题 20 分
%% typeId: 6
%% type: 计算题
%% chapter: 磁场
%% point: 带电粒子在磁场中的运动
%% method: 无
%% score: 20
%% degree: 350
%% duplicateId: 0
%% body:
如图，在 $y>0$ 的区域存在方向沿 $y$ 轴负方向的匀强电场，场强大小为 $E$，在 $y<0$ 的区域存在方向垂直于 $xOy$ 平面向外的匀强磁场。一个氕核 \ce{^1_1 H} 和一个氘核 \ce{^2_1 H} 先后从 $y$ 轴上 $y=h$ 点以相同的动能射出，速度方向沿 $x$ 轴正方向。已知 \ce{^1_1 H} 进入磁场时，速度方向与 $x$ 轴正方向的夹角为 $60\Udu$，并从坐标原点 $O$ 处第一次射出磁场。\ce{^1_1 H} 的质量为 $m$，电荷量为 $q$，不计重力。求：
\begin{enumerate}
	\item
	\ce{^1_1 H} 第一次进入磁场的位置到原点 $O$ 的距离；
	\item
	磁场的磁感应强度大小；
	\item
	\ce{^2_1 H} 第一次离开磁场的位置到原点 $O$ 的距离。
\end{enumerate}

\onepicture[align=right]{figs/25.png}

%% answer:
\jdanswer{
\begin{enumerate}
	\item
	$s_1 = \frac{2\sqrt{3}}{3}h$

	\item
	$B = \sqrt{\frac{6mE}{qh}}$

	\item
	$s_2' - s_2 = \frac{2\sqrt{3}}{3}(\sqrt{2}-1)h$

\end{enumerate}
}
%% memo:
\memoanswer{
（1）\ce{^1_1 H} 在电场中做类平抛运动，在磁场中做圆周运动，运动轨迹如图所示。设 \ce{^1_1 H} 在电场中的加速度大小为 $a_1$，初速度大小为 $v_1$，它在电场中的运动时间为 $t_1$，第一次进入磁场的位置到原点 $O$ 的距离为 $s_1$。由运动学公式得

$s_1=v_1t_1$\ccnum{①}，

$h=\frac{1}{2}a_1t_1^2$\ccnum{②}，

\onepicture[w=3.4cm]{figs/25_an.jpg}

由题给条件，\ce{^1_1 H} 进入磁场时速度的方向与 $x$ 轴正方向夹角 $\theta_1=60\Udu$，\ce{^1_1 H} 进入磁场时速度的 $y$ 分量的大小为

$a_1t_1=v_1\tan\theta_1$\ccnum{③}，

联立以上各式得 $s_1=\frac{2\sqrt{3}}{3}h$\ccnum{④}。

（2）\ce{^1_1 H} 在电场中运动时，由牛顿第二定律得

$qE=ma_1$\ccnum{⑤}，

设 \ce{^1_1 H} 进入磁场时速度的大小为 $v_1'$，由速度合成法则得

$v_1'=\sqrt{v_1^2+(a_1t_1)^2}$\ccnum{⑥}，

设磁感应强度大小为 $B$，\ce{^1_1 H} 在磁场中运动的圆轨道半径为 $R_1$，由洛伦兹力公式和牛顿第二定律得

$qv_1'B=\frac{mv_1'^2}{R_1}$\ccnum{⑦}，

由几何关系得

$s_1=2R_1\sin\theta_1$\ccnum{⑧}，

联立以上各式得 $B=\sqrt{\frac{6mE}{qh}}$\ccnum{⑨}。

（3）设 \ce{^2_1 H} 在电场中沿 $x$ 轴正方向射出的速度大小为 $v_2$，在电场中的加速度大小为 $a_2$，由题给条件得

$\frac{1}{2}(2m)v_2^2=\frac{1}{2}mv_1^2$\ccnum{⑩}，

由牛顿第二定律得

$qE=2ma_2$\ccnum{⑪}，

设 \ce{^2_1 H} 第一次射入磁场时的速度大小为 $v_2'$，速度的方向与 $x$ 轴正方向夹角为 $\theta_2$，入射点到原点的距离为 $s_2$，在电场中运动的时间为 $t_2$，由运动学公式得

$s_2=v_2t_2$\ccnum{⑫}，

$h=\frac{1}{2}a_2t_2^2$\ccnum{⑬}，

$v_2'=\sqrt{v_2^2+(a_2t_2)^2}$\ccnum{⑭}，

$\sin\theta_2=\frac{a_2t_2}{v_2}$\ccnum{⑮}，

$s_2=s_1$，$\theta_2=\theta_1$，$v_2'=\frac{\sqrt{2}}{2}v_1'$\ccnum{⑯}。

设 \ce{^2_1 H} 在磁场中做圆周运动的半径为 $R_2$，由\ccnum{⑦}\ccnum{⑯}式及粒子在匀强磁场中做圆周运动的半径公式得

$R_2=\frac{(2m)v_2'}{qB}=\sqrt{2}R_1$\ccnum{⑰}，

所以出射点在原点左侧。设 \ce{^2_1 H} 进入磁场的入射点到第一次离开磁场的出射点的距离为 $s_2'$，由几何关系得

$s_2'=2R_2\sin\theta_2$\ccnum{⑱}，

联立\ccnum{④}\ccnum{⑧}\ccnum{⑯}\ccnum{⑰}\ccnum{⑱}式得，\ce{^2_1 H} 第一次离开磁场时的位置到原点 $O$ 的距离为

$s_2'-s_2=\frac{2\sqrt{3}}{3}(\sqrt{2}-1)h$。
}

\item
%% number: 33
%% paperName: 2018年全国理综Ⅰ第 33 题 10 分
%% typeId: 6
%% type: 计算题
%% chapter: 气体性质
%% point: 气体连接体
%% method: 无
%% score: 10
%% degree: 950
%% duplicateId: 0
%% body:
（1）如图，一定质量的理想气体从状态 a 开始，经历过程①、②、③、④到达状态 e，对此气体，下列说法正确的是\xzanswer{BDE}

\onepicture[w=6.0cm]{figs/33a.png}

\fivechoices[answer=BDE]
{过程①中气体的压强逐渐减小}
{过程②中气体对外界做正功}
{过程④中气体从外界吸收了热量}
{状态 c、d 的内能相等}
{状态 d 的压强比状态 b 的压强小}

（2）如图，容积为 $V$ 的汽缸由导热材料制成，面积为 $S$ 的活塞将汽缸分成容积相等的上下两部分，汽缸上部通过细管与装有某种液体的容器相连，细管上有一阀门 K。开始时，K 关闭，汽缸内上下两部分气体的压强均为 $p_0$。现将 K 打开，容器内的液体缓慢地流入汽缸，当流入的液体体积为 $\frac{V}{8}$ 时，将 K 关闭，活塞平衡时其下方气体的体积减小了 $\frac{V}{6}$，不计活塞的质量和体积，外界温度保持不变，重力加速度大小为 $g$。求流入汽缸内液体的质量。

\onepicture[h=4.2cm, align=right]{figs/33b.png}

%% answer:
\jdanswer{
$m = \frac{15p_0S}{26g}$
}
%% memo:
\memoanswer{
（1）过程①气体发生等容变化，温度升高，由 $\frac{pV}{T}=C$ 可知气体压强增大，故 A 错误；过程②气体体积增大，气体对外做正功，故 B 正确；过程④气体发生等容变化，气体对外不做功，温度降低，内能减小，由 $\Delta U=Q+W$ 可知气体对外放热，故 C 错误；状态 c、d 的温度相同，气体内能相等，故 D 正确；由 $\frac{pV}{T}=C$ 可得 $T=\frac{p}{c}V$，在 $T-V$ 图像中，坐标点与坐标原点的连线的斜率 $k=\frac{p}{c}$，如图所示，所以状态 d 的压强比状态 b 的压强小，故 E 正确。

\onepicture[w=4.5cm]{figs/33_an.jpg}

（2）设活塞再次平衡后，活塞上方气体的体积为 $V_1$，压强为 $p_1$；下方气体的体积为 $V_2$，压强为 $p_2$。在活塞下移的过程中，活塞上、下方气体的温度均保持不变，由玻意耳定律得

$p_0\frac{V}{2}=p_1V_1$\ccnum{①}，

$p_0\frac{V}{2}=p_2V_2$\ccnum{②}，

由已知条件得

$V_1=\frac{V}{2}+\frac{V}{6}-\frac{V}{8}=\frac{13}{24}V$\ccnum{③}，

$V_2=\frac{V}{2}-\frac{V}{6}=\frac{V}{3}$\ccnum{④}，

设活塞上方液体的质量为 $m$，由力的平衡条件得

$p_2S=p_1S+mg$\ccnum{⑤}，

联立以上各式得 $m=\frac{15p_0S}{26g}$\ccnum{⑥}。
}

\item
%% number: 34
%% paperName: 2018年全国理综Ⅰ第 34 题 10 分
%% typeId: 6
%% type: 计算题
%% chapter: 机械波
%% point: 波与振动图象
%% method: 无
%% score: 10
%% degree: 550
%% duplicateId: 0
%% body:
（1）如图，$\triangle ABC$ 为一玻璃三棱镜的横截面，$\angle A = 30\Udu$，一束红光垂直 AB 边射入，从 AC 边上的 D 点射出，其折射角为 $60\Udu$，则玻璃对红光的折射率为\tkanswer[1.4cm]{$\sqrt{3}$}。若改用蓝光沿同一路径入射，则光线在 D 点射出时的折射角\tkanswer[1.6cm]{大于}（“小于”“等于”或“大于”）$60\Udu$。

\onepicture[h=3.4cm, align=right]{figs/34a.png}

（2）一列简谐横波在 $t = \frac{1}{3}\Us$ 时的波形图如图（a）所示，P、Q 是介质中的两个质点，图（b）是质点 Q 的振动图像。求：
\begin{enumerate}
	\item
	波速及波的传播方向；
	\item
	质点 Q 的平衡位置的 $x$ 坐标。
\end{enumerate}

\twopicture[labelA=2018全国理综Ⅰ34b, labelB=2018全国理综Ⅰ34c, numA={（a）}, numB={（b）}, hA=2.6cm, hB=2.6cm, align=right]
{figs/34b.png}
{figs/34c.png}

%% answer:
\jdanswer{
\begin{enumerate}
	\item
	$v = 18\Ucms$，波沿 $x$ 轴负方向传播

	\item
	$x_Q = 9\Ucm$

\end{enumerate}
}
%% memo:
\memoanswer{
（1）由几何知识可得红光从玻璃射入空气中的入射角为 $i = 30\Udu$，折射角为 $r = 60\Udu$，由光的折射定律可得 $n = \frac{\sin r}{\sin i} = \sqrt{3}$。若改用蓝光入射，在玻璃中入射角不变，由于蓝光的频率比红光的频率大，所以玻璃对蓝光的折射率比对红光的折射率大，由光的折射定律可得光线在 D 点射出时的折射角大于 $60\Udu$。

（2）（i）由图（a）可以看出，该波的波长为

$\lambda = 36\Ucm$\ccnum{①}，

由图（b）可以看出，周期为

$T = 2\Us$\ccnum{②}，

波速为 $v = \frac{\lambda}{T} = 18\Ucms$\ccnum{③}，

由图（b）知，当 $t = \frac{1}{3}\Us$ 时，Q 点向上运动，结合图（a）可得，波沿 $x$ 轴负方向传播。

（ii）设质点 P、Q 平衡位置的 $x$ 坐标分别为 $x_P$、$x_Q$。由图（a）知，$x = 0$ 处 $y = -\frac{A}{2} = A\sin(-30\Udu)$，因此

$x_P = \frac{30\Udu}{360\Udu}\lambda = 3\Ucm$\ccnum{④}，

由图（b）知，在 $t = 0$ 时 Q 点处于平衡位置，经 $\Delta t = \frac{1}{3}\Us$，其振动状态向 $x$ 轴负方向传播至 P 点处，由此及\ccnum{③}式有

$x_Q - x_P = v\Delta t = 6\Ucm$\ccnum{⑤}，

由\ccnum{④}\ccnum{⑤}式得，质点 Q 的平衡位置的 $x$ 坐标为

$x_Q = 9\Ucm$\ccnum{⑥}。
}

\end{enumerate}

\end{document}
